\begin{ex}
Trong các số: $45, 54, 49, 51$, số lớn nhất là số nào?
\shortans{54}
\loigiai{
So sánh chữ số hàng chục: $5 > 4$. Hai số có hàng chục là 5 là 54 và 51, trong đó $54 > 51$. Vậy số lớn nhất là $54$.
\textbf{Đáp số:} $54$.
}
\end{ex}

\begin{ex}
Trong các số: $78, 87, 69, 90$, số bé nhất là số nào?
\shortans{69}
\loigiai{
Chữ số hàng chục nhỏ nhất là 6 (trong số 69). Vậy số bé nhất là $69$.
\textbf{Đáp số:} $69$.
}
\end{ex}

\begin{ex}
Tìm số tự nhiên $x$ thỏa mãn: $38 < x < 40$.
\shortans{39}
\loigiai{
Số tự nhiên nằm giữa 38 và 40 là số $39$.
\textbf{Đáp số:} $39$.
}
\end{ex}

\begin{ex}
Cho các số: $23, 67, 45, 89, 12$. Sắp xếp theo thứ tự từ lớn đến bé, số đứng ở vị trí thứ ba là số nào?
\shortans{45}
\loigiai{
Thứ tự từ lớn đến bé: $89 > 67 > 45 > 23 > 12$. Số đứng thứ ba là số $45$.
\textbf{Đáp số:} $45$.
}
\end{ex}

\begin{ex}
Cho các số: $72, 35, 56, 91, 48$. Sắp xếp theo thứ tự từ bé đến lớn, số đứng ở vị trí thứ hai là số nào?
\shortans{48}
\loigiai{
Thứ tự từ bé đến lớn: $35 < 48 < 56 < 72 < 91$. Số đứng thứ hai là số $48$.
\textbf{Đáp số:} $48$.
}
\end{ex}

\begin{ex}
Tìm số tròn chục lớn hơn 55 và bé hơn 65.
\shortans{60}
\loigiai{
Trong khoảng từ 55 đến 65 chỉ có duy nhất số tròn chục $60$.
\textbf{Đáp số:} $60$.
}
\end{ex}

\begin{ex}
Tìm số tròn chục lớn hơn 72 và bé hơn 88.
\shortans{80}
\loigiai{
Số tròn chục thỏa mãn là số $80$.
\textbf{Đáp số:} $80$.
}
\end{ex}

\begin{ex}
Tìm số lớn nhất có hai chữ số mà chữ số hàng chục là 7.
\shortans{79}
\loigiai{
Để số lớn nhất, chữ số hàng đơn vị phải lớn nhất là 9. Vậy số đó là $79$.
\textbf{Đáp số:} $79$.
}
\end{ex}

\begin{ex}
Tìm số bé nhất có hai chữ số mà chữ số hàng chục là 5.
\shortans{50}
\loigiai{
Để số bé nhất, chữ số hàng đơn vị phải bé nhất là 0. Vậy số đó là $50$.
\textbf{Đáp số:} $50$.
}
\end{ex}

\begin{ex}
Có bao nhiêu số có hai chữ số lớn hơn 85 và bé hơn 90?
\shortans{4}
\loigiai{
Các số đó là: $86, 87, 88, 89$. Có tất cả $4$ số.
\textbf{Đáp số:} $4$.
}
\end{ex}

\begin{ex}
Có bao nhiêu số tròn chục lớn hơn 25 và bé hơn 75?
\shortans{5}
\loigiai{
Các số tròn chục đó là: $30, 40, 50, 60, 70$. Có tất cả $5$ số.
\textbf{Đáp số:} $5$.
}
\end{ex}

\begin{ex}
Cho ba số: $a = 64, b = 46, c = 66$. Sắp xếp ba số theo thứ tự từ bé đến lớn, số đứng ở đầu tiên là số nào?
\shortans{46}
\loigiai{
So sánh: $46 < 64 < 66$. Số bé nhất đứng đầu tiên là $46$.
\textbf{Đáp số:} $46$.
}
\end{ex}

\begin{ex}
Lấy số lớn nhất trong các số: $35, 82, 53, 28$ trừ đi số bé nhất trong các số đó được kết quả bằng bao nhiêu?
\shortans{54}
\loigiai{
Số lớn nhất là 82, số bé nhất là 28: $82 - 28 = 54$.
\textbf{Đáp số:} $54$.
}
\end{ex}

\begin{ex}
Tìm số tự nhiên $y$ tròn chục sao cho: $40 < y < 60$.
\shortans{50}
\loigiai{
Số tròn chục nằm giữa 40 và 60 là số $50$.
\textbf{Đáp số:} $50$.
}
\end{ex}

\begin{ex}
Trong một đợt kiểm tra, bạn An được 85 điểm, bạn Bình được 90 điểm, bạn Cúc được 88 điểm. Bạn nào đạt điểm cao nhất? (An: 1, Bình: 2, Cúc: 3).
\shortans{2}
\loigiai{
So sánh điểm: $90 > 88 > 85$. Bạn Bình đạt điểm cao nhất (ghi số 2).
\textbf{Đáp số:} $2$.
}
\end{ex}

\begin{ex}
Tìm số có hai chữ số giống nhau lớn hơn 70 và bé hơn 80.
\shortans{77}
\loigiai{
Số có hai chữ số giống nhau trong khoảng này là số $77$.
\textbf{Đáp số:} $77$.
}
\end{ex}

\begin{ex}
Tìm số có hai chữ số giống nhau lớn hơn 80 và bé hơn 95.
\shortans{88}
\loigiai{
Số có hai chữ số giống nhau thỏa mãn là số $88$.
\textbf{Đáp số:} $88$.
}
\end{ex}

\begin{ex}
Hiệu giữa số lớn nhất có hai chữ số và số lớn nhất có một chữ số là bao nhiêu?
\shortans{90}
\loigiai{
Số lớn nhất có hai chữ số là 99, số lớn nhất có một chữ số là 9: $99 - 9 = 90$.
\textbf{Đáp số:} $90$.
}
\end{ex}

\begin{ex}
Hiệu giữa số bé nhất có hai chữ số và số bé nhất có một chữ số khác 0 là bao nhiêu?
\shortans{9}
\loigiai{
Số bé nhất có hai chữ số là 10, số bé nhất có một chữ số khác 0 là 1: $10 - 1 = 9$.
\textbf{Đáp số:} $9$.
}
\end{ex}

\begin{ex}
Tìm số tự nhiên $x$ có hai chữ số biết: $x$ lớn hơn 97 và bé hơn 99.
\shortans{98}
\loigiai{
Số tự nhiên nằm giữa 97 và 99 là số $98$.
\textbf{Đáp số:} $98$.
}
\end{ex}

\begin{ex}
Có bao nhiêu số có hai chữ số có chữ số tận cùng là 0?
\shortans{9}
\loigiai{
Đó là các số tròn chục: $10, 20, 30, \dots, 90$. Có tất cả $9$ số.
\textbf{Đáp số:} $9$.
}
\end{ex}

\begin{ex}
Tìm số tròn chục liền trước của số 100.
\shortans{90}
\loigiai{
Các số tròn chục: $80, 90, 100$. Số tròn chục liền trước 100 là $90$.
\textbf{Đáp số:} $90$.
}
\end{ex}

\begin{ex}
Hai bạn Mai và Lan sưu tầm tem. Mai có 52 con tem, Lan có 48 con tem. Hỏi Mai nhiều hơn Lan bao nhiêu con tem?
\shortans{4}
\loigiai{
Số tem Mai nhiều hơn Lan là: $52 - 48 = 4$ (con tem).
\textbf{Đáp số:} $4$.
}
\end{ex}

\begin{ex}
Bác An thu hoạch được 65 bắp ngô, bác Ba thu hoạch được 70 bắp ngô. Bác Ba thu hoạch nhiều hơn bác An bao nhiêu bắp ngô?
\shortans{5}
\loigiai{
Số bắp ngô bác Ba nhiều hơn bác An là: $70 - 65 = 5$ (bắp ngô).
\textbf{Đáp số:} $5$.
}
\end{ex}

\begin{ex}
Cho dãy số: $15, 20, 25, 30, \dots, 40$. Số thích hợp điền vào chỗ chấm là số nào?
\shortans{35}
\loigiai{
Dãy số cách đều 5 đơn vị: $30 + 5 = 35$.
\textbf{Đáp số:} $35$.
}
\end{ex}
